\subsection{Interior Regularity}\label{sec:elliptic-interior-regularity}

In Section~\ref{sec:dirichlet-existence-and-uniqueness}, we have shown the existence of weak solutions $u\in H^k(\Omega)$ of the Dirichlet problem for elliptic operators of order $2k$. We now wish to show that under suitable hypotheses on the smoothness of the coefficients $a_{\sigma\gamma}$, the forcing function $f$ and the boundary of $\Omega$, our weak solution is, in fact, a strong solution or a classical solution. In order to give some idea of how we plan to go about this, we make a couple of formal calculations.

For our first calculation let us assume that $\Omega$ has a smooth boundary $\partial\Omega$ with unit outward normal $\bm{\eta}=(\eta_1,\ldots,\eta_n)$ and that $u$ is a classical solution of
\begin{equation}\label{eq:interior-regularity-formal-poisson-equation}
-\Delta u=f
\end{equation}
in $\Omega$, and
\begin{equation}\label{eq:interior-regularity-formal-dirichlet-boundary}
u=0
\end{equation}
on $\partial\Omega$. Our goal is to show that (weak) solutions of elliptic problems such as the one above are actually in a ``better\textquotesingle\textquotesingle\ space than $H_0^1(\Omega)$. In order to prepare for this, we will now estimate the $L^2(\Omega)$ norm of the matrix of second partials of $u$ in terms of the $H^1(\Omega)$ norm. Since this is simply a formal calculation, we will proceed as if we already know that $u$ is as smooth as we like.

\[\begin{aligned}
\int_\Omega|\Delta u|^2\,\mathrm{d}\bm{x}
&=\sum_{i,j=1}^n\int_\Omega u_{x_ix_i}u_{x_jx_j}\,\mathrm{d}\bm{x}\\
&=-\sum_{i,j=1}^n\int_\Omega u_{x_i}u_{x_jx_ix_j}\,\mathrm{d}\bm{x}\\
&\quad+\sum_{i,j=1}^n\int_{\partial\Omega}u_{x_i}u_{x_jx_j}\eta_i\,\mathrm{d}S\\
&=\sum_{i,j=1}^n\int_\Omega u_{x_ix_j}u_{x_jx_i}\,\mathrm{d}\bm{x}\\
&\quad+\sum_{i,j=1}^n\int_{\partial\Omega}u_{x_i}u_{x_jx_j}\eta_i-u_{x_i}u_{x_ix_j}\eta_j\,\mathrm{d}S.
\end{aligned}\]

We also have
\begin{equation}\label{eq:interior-regularity-formal-laplacian-data-identity}
\int_\Omega|\Delta u|^2\,\mathrm{d}\bm{x}=\int_\Omega|f|^2\,\mathrm{d}\bm{x}.
\end{equation}
Combining these two results gives us
\begin{equation}\label{eq:interior-regularity-formal-hessian-boundary-estimate}
\int_\Omega|\nabla^2u|^2\,\mathrm{d}\bm{x}\leq\int_\Omega|f|^2\,\mathrm{d}\bm{x}+|\text{boundary terms}|.
\end{equation}
Thus, if we had some additional information on the boundary terms, we could derive an a priori estimate on the $H^2(\Omega)$ norm of a solution $u$ in terms of the data $f$.
